\section{Related work}

The closest asymptotic efficiency framework is that of \citet[Theorems 3.3 and 3.5, Remark 3.6, and Section 4]{Steinberger2024}. For parametric models differentiable in quadratic mean, his sequential theory establishes local asymptotic mixed normality along subsequences with convergent conditional information, under countably generated input and output sigma-fields. The convolution theorem additionally imposes positive-definiteness conditions on the conditional information matrices and their limiting accumulation points. His scalar optimization maximizes output Fisher information, and the private two-step procedure uses a preliminary estimate to select an asymptotically optimal mechanism under the stated regularity conditions. Remark 3.6 explicitly identifies scalar functionals of inverse information as objectives for multiparameter models. % [Efficiency in local differential privacy](https://arxiv.org/html/2301.10600v3)
Our analysis treats a treatment-effect contrast in a two-parameter \leanref{S-1}{randomized binary trial} with binary outcomes and known interior assignment probability, as specified in \cref{obj:def:binary-trial-model}. At interior arm means and a fixed positive privacy budget, \cref{obj:thm:finite-oracle,obj:thm:sequential-minimax} connect finite channel optimization with the exact efficiency bound for the stated class of regular procedures using arbitrary measurable sequential output spaces. The scalar target therefore retains a nuisance parameter throughout the optimization.

The finite reduction is closely related to the staircase mechanisms of \citet[Theorems 2 and 4]{KairouzOhViswanath2016}. Their utility class consists of sums over output columns of a sublinear function, meaning a positively homogeneous, subadditive function of the column. For this class, they establish an optimal staircase mechanism with output alphabet no larger than the input alphabet and express the optimum as a finite linear program. % [Extremal mechanisms for local differential privacy](https://jmlr.org/papers/volume17/15-135/15-135.pdf)
The present \leanref{S-3}{fourteen-pattern staircase program}, characterized in \cref{obj:synth_4,obj:def:staircase-saddle,obj:thm:finite-oracle}, combines this finite channel geometry with nuisance profiling. It also parallels contrast-specific optimal-design geometry: \citet[Sections 2--4]{Elfving1952} represents optimal designs through a convex set generated by signed regression vectors, and \citet[Chapter 2]{Pukelsheim2006} develops the corresponding information-matrix framework. For a fixed parameter direction, directional Fisher information here is additive over staircase outputs and linear in their weights. Contrast information minimizes that directional information over directions normalized to have unit treatment-effect derivative. This minimization makes the objective concave in the staircase weights and gives the optimization its saddle structure. The relevant support geometry is consequently governed jointly by channel normalization and profiling. Under the conditions stated in the support results, \cref{obj:thm:five-output-upper-bound} supplies a five-output upper bound, while \cref{obj:thm:five-output-region} identifies five-output optimality through its certificate and score-separation conditions. Comparison by randomization provides the broader experiment-ordering context \citep[pp. 265--272]{Blackwell1953}; the results characterize contrast-specific efficiency within this framework.

The nuisance projection and local risk formulation follow the local-asymptotic tradition of \citet[Chapter 6]{LeCam2000} and \citet[Chapters 7--8 and 25]{VanDerVaart1998}. The sequential lower bound uses the Bayesian information-inequality lineage of \citet{VanTrees2001} and \citet[pp. 59--79]{Gill1995}. These references provide intellectual antecedents; the formal verification contract records no external theorem as a Lean proof dependency.

The direct causal benchmark is \citet[Sections 3.1 and 3.3, Lemma 4, Theorem 7, and Appendix A.5]{OhnishiAwan2025}. Their randomized-experiment model permits potential outcomes in the unit interval. Lemma 4 gives a worst-case squared-error lower bound for privacy budgets between zero and one. In the known-assignment custom scenario, each subject forms the inverse-probability-weighted observed outcome and adds independent Laplace noise calibrated to the sum of the reciprocal treatment and control assignment probabilities. The sample mean of these releases is unbiased and consistent; Theorem 7 supplies asymptotic normality, a variance estimate, confidence intervals, and the privacy-dependent convergence rate. These statements provide both a rate benchmark and a concrete inferential procedure. % [Locally private causal inference for randomized experiments](https://jmlr.org/papers/volume26/23-1401/23-1401.pdf)
For binary outcomes, our comparison uses precisely this release and sample-mean estimator, specified in \cref{obj:def:oa-release}. The comparison in \cref{obj:thm:laplace-comparison} evaluates its asymptotic variance against the optimized contrast bound at the stated trial parameters and privacy budget.

These exact variance comparisons complement the literature on statistical rates under local privacy. In the bounded scalar-mean problem, \citet[Section 3.2.1]{DuchiJordanWainwright2018} obtain squared-error order \(\min\{1,(n\varepsilon^2)^{-1}\}\), up to numerical constants, in the high-privacy regime \(0<\varepsilon\le1\). \Citet[Sections 3.1 and 3.3, Lemma 4]{OhnishiAwan2025} obtain the same benchmark order for their bounded-outcome randomized-experiment model. This rate describes the joint sample-size/privacy scaling; the fixed-privacy efficiency constant derived here resolves the model- and design-specific numerical factor that rate notation suppresses. \Citet{Rohde2020} characterize the difficulty of functional estimation through total-variation moduli, including rate-optimal binary releases for linear functionals over convex models. Instance-specific local minimax complexity is developed further by \citet{Duchi2024}. % [Minimax optimal procedures for locally private estimation](https://arxiv.org/abs/1604.02390), [Geometrizing rates of convergence under local differential privacy constraints](https://arxiv.org/abs/1805.01422), [The right complexity measure in locally private estimation](https://arxiv.org/abs/1806.05756)
Information-based approaches also include the Fisher-information contraction inequalities of \citet{BarnesChenOzgur2020} and the communication-based estimation lower bounds of \citet{Duchi2019}. % [Fisher information under local differential privacy](https://arxiv.org/abs/2005.10783), [Lower bounds for locally private estimation via communication complexity](https://proceedings.mlr.press/v99/duchi19a.html)
Here the finite binary-trial experiment permits an exact optimization of the asymptotic contrast variance at a fixed privacy budget.

Exact scalar mechanisms provide a further point of comparison. For a Gaussian mean with known unit variance, \citet{KalininSteinberger2025} identify randomized response applied to the centered sign as an information-maximizing mechanism in their high-privacy regime, and attain efficiency through a private preliminary estimate. % [Efficient estimation of a Gaussian mean with local differential privacy](https://proceedings.mlr.press/v258/nikita25a.html)
The private-pilot construction in \cref{obj:def:pilot-estimator,obj:thm:pilot-attainment} uses the same broad principle of learning the parameter before selecting the main-sample release. Its channel selection solves the nuisance-profiled contrast problem, with pilot size diverging and occupying a vanishing fraction of the sample under \cref{obj:ass:pilot-diverges,obj:ass:pilot-sublinear}.

Recent work further develops extremal-channel optimization. \citet{Amorino2025} study factorization through extremal mechanisms and Fisher-information maximization in discrete and continuous settings, while \citet{Amorino2026} analyze symmetry and asymmetry in staircase information across privacy regimes. % [Factorization by extremal privacy mechanisms](https://arxiv.org/abs/2507.21769), [On the role of symmetry for staircase mechanisms](https://arxiv.org/abs/2603.27572)
For binary hypothesis testing between two known finite-alphabet distributions, \citet{GhaziNasserCalmonIssa2026} optimize output \(f\)-divergences through likelihood-ratio ordering, contiguous partitions, and randomized response. % [Sort, partition, randomize](https://arxiv.org/abs/2606.07443)
These objectives illuminate the role of staircase structure. Our contrast objective additionally accounts for variation in both arm means, and its support certificates characterize the resulting release cardinality.

Finally, causal efficiency connects the analysis to \citet{Hahn1998,Hirano2003}, whose work studies efficient average treatment-effect estimation under unconfoundedness. % [Efficient estimation of average treatment effects using the estimated propensity score](https://scholar.harvard.edu/files/imbens/files/efficient_estimation_of_average_treatment_effects_using_the_estimated_propensity_score.pdf)
In the private setting, \citet{KormosVanderVaart2026} construct suitable identifying mechanisms that transfer rate-double-robust and semiparametric inferential properties to noisy observations, including for causal targets. % [Private rate-double-robust inference](https://arxiv.org/abs/2606.20427)
The present binary randomized-trial model makes the release mechanism itself an object of efficiency optimization: the channel, nuisance adjustment, and attainable treatment-effect variance are characterized together under the stated sequential privacy conditions.
